\textbf{证据确立了什么。}源代码多样性、执行机制和单次正确性差异是不同的测量。Phase-I 模式未在调整后的 Phase-II 设置中复现，但多项协议变化使其原因无法识别。原门控的修正分析未显示明确的提升空间改善；该门控还排除了允许的提示策略，因此无法识别单独验证的效应。

\textbf{尚未识别什么。}同代码重跑可能发生分歧并产生 oracle 增益。结果翻转率测量不稳定性，不是处理对比的标准误，也不能定义其数值噪声地板。在将提升空间归因于稳定新增能力前，需要匹配成本的克隆对照、独立重复验证，以及经过验证的跨机制类型校准（附录~\ref{app:gatecontrols}）。

\textbf{局限。}主比较使用一个文本转 SQL 目标和六个固定构建器，每个构建器三个生成种子。Phase I 属于发现证据，存在已披露的任务重叠、代理评分器和查看轨迹后设计的人工控制。审后新统计分析未增加独立执行。K-matched 与 R2 已共用修正数据和重采样链条，W1/W3 也已共用 canonical 数据，但仍是描述性分析；选择不确定性与独立效用验证尚未完成。已完成的单种子 MATH-500 自由组探查（附录~\ref{app:crossdomain}）将观测结果比较扩展至第二领域，但不提供克隆对照、修复门控的析因实验或稳定互补性的独立检验。
